\section{Introduction: concentration and power channels}\label{pr:section}

Diffraction provides familiar scales for focusing and for the number of
spatial patterns communicated through an aperture. Inhomogeneous dielectric
media also support strong near fields and resonant energy storage. Motivated
by OQ14 of Miller and Monticone~\cite[Section 5.6]{millerMonticone2026},
we study two related questions: how many localized electric-field maxima
can a single field produce, and how many independent inputs can carry power
through a prescribed section? Both questions require explicit choices of
normalization, geometric resources, and admissibility conditions.

For the \emph{hotspot} question, we seek finitely many disjoint observation
regions and one field with a large, localized electric-field maximum in each.
A local singularity at a dielectric corner suggests a mechanism,
but a packing theorem also requires a globally admissible field that excites
all regions and retains the peaks after rounding. We distinguish driven
fields, normalized by incoming power, from outgoing resonant modes, for
which a different normalization is required. No fixed fraction of the
total energy is required in each peak.

For the \emph{channel} question, fix a square $D$ in a plane and consider
a linear space of incident fields. Every field in the space must satisfy
three tests: a lower bound on power in selected outputs, a lower bound on
signed power flux through $D$, and an upper bound on the integral of absolute
flux over the rest of the plane. Requiring these inequalities for each basis
vector alone would not control their coherent superpositions. The square is
a measurement section and need not be the only topological route between
input and output.

Figure~\ref{fig:channel-geometry} shows the measurement geometry and the
distinction between flux through the square and flux elsewhere in its plane.
\begin{figure}[htbp]
\centering
\setlength{\unitlength}{1mm}
\begin{picture}(150,65)
  \put(2,45){\makebox(38,6)[l]{\small View of the plane $z=0$}}
  \put(3,7){\framebox(66,37){}}
  \put(26,16){\framebox(20,20){\large $D$}}
  \put(28,10){\makebox(16,5){\scriptsize fixed square}}
  \put(6,37){\makebox(19,5)[l]{\scriptsize $\Pi\setminus D$}}
  \put(4,1){\makebox(65,5)[l]{\scriptsize plane continues beyond frame}}
  \put(78,45){\makebox(70,6)[l]{\small Power accounting}}
  \put(79,37){\vector(1,0){12}}
  \put(76,41){\makebox(21,4){\scriptsize input $a$}}
  \put(92,31){\framebox(26,14){\shortstack{\scriptsize dielectric\\[-1pt]\scriptsize device}}}
  \put(119,37){\vector(1,0){12}}
  \put(131,33){\makebox(18,8)[l]{\shortstack[l]{\scriptsize selected\\[-1pt]\scriptsize output $Ta$}}}
  \put(79,19){\makebox(68,7)[l]{\scriptsize On $D$: signed forward flux}}
  \put(79,11){\makebox(68,7)[l]{\scriptsize Outside $D$: absolute flux}}
  \put(79,3){\makebox(68,7)[l]{\scriptsize Inequalities hold for every $a\in W$}}
\end{picture}
\caption{Schematic of the channel definition. The left panel is a view of
the whole measurement plane, with a fixed square $D$; the outer frame only
marks the drawing boundary. The right panel indicates input--output power
accounting, not a prescribed one-way device geometry. Passage uses signed
flux through $D$, whereas bypass measures the absolute flux on the rest of
the plane. In the counterexample the selected output is the complete
outgoing scattering map, including reflection.}
\label{fig:channel-geometry}
\end{figure}


Across each device family, the driving frequency or the real part of the
resonance frequency, as appropriate, and the material bounds are fixed.
For channels, the square and the thresholds are also fixed. The surrounding
device may grow; the formulation imposes no
positive lower bound on feature size and no uniform upper bound on quality
factor, number of interfaces, or matching complexity. A bound depending on
one of these additional resources addresses a more restricted question.
The precise admissibility and power conventions follow below.

\subsection{Main results and their meaning}
The following summary is informal; the linked statements give the exact
hypotheses and conclusions.
\begin{itemize}
\item \textbf{Driven hotspots.}  Theorem~\ref{hs:packing} gives
  simultaneous peaks in a two-dimensional invariant geometry, with incoming
  power normalized per unit invariant length.  Theorem~\ref{hs:three-packing}
  gives a finite three-dimensional construction driven by one unit-power
  illumination.  In that theorem the observation neighborhoods project into
  a prescribed planar patch; their maxima need not lie on the plane.
\item \textbf{Modal hotspots.}  Theorem~\ref{ms:modal-packing} constructs
  arbitrarily many localized peaks in one normalized outgoing mode of a
  finite smooth dielectric.  The real part of its resonance frequency, its
  quality-factor target and common material bounds can be prescribed. The
  norm is a reciprocal resonance-residue normalization, not an all-space
  positive-energy norm;
  the modal-volume conclusions are stated separately in
  Corollary~\ref{ms:modal-volume}. Equivalence to other possible
  interpretations of the original modal-volume question is not established.
\item \textbf{Unrestricted channels.}  Theorem~\ref{uc:counterexample}
  constructs an $N$-dimensional counted input subspace for every finite
  $N$, while the frequency, material interval, square and thresholds remain
  fixed.  It applies when $0<\eta<\delta$, with the complete outgoing
  scattering map selected.  Thus the proposed uniform finite count, as
  quantified below over these allowed parameters, is false.  This does not
  settle $\delta\leq\eta$, a specified receiving-only output, or fixed total
  device support. These results do not settle the original broad OQ14 for
  every possible interpretation of its undefined thresholds and outputs.
\end{itemize}
For intuition, the channel construction uses a distant radial dielectric
stack to amplify independent local field patterns.  Small dielectric
inclusions turn those patterns into local flux circulation near the square's
edge.  The proof must control both the full self-consistent field and the
absolute flux over the \emph{whole} plane, uniformly over coherent inputs;
a large local field by itself would not establish the conclusion.

\subsection{How to read the argument}
Sections~\ref{hs:section} and~\ref{hs:three-section} develop the driven
hotspot constructions.  Sections~\ref{hs:modal-next-section}
and~\ref{ms:section} treat excitation and normalization of a single
resonant mode.  Section~\ref{ch:section} gives channel bounds that retain an
explicit field resource; the intervening channel sections examine the
passage, control and physical-construction issues.  The unrestricted
counterexample is stated and proved in Section~\ref{uc:section}.
Appendix~\ref{ver:section} separates mathematical results, formal checks
and numerical diagnostics.  A reader interested in the central conclusion
can read this introduction and the theorem statements first, then return to
the supporting estimates.

\subsection{Precise targets and conventions}
The source question does not fix every normalization, threshold, geometric
resource, or interpretation of a mode.  The definitions below make our
choices explicit; they are not quotations of an already quantified theorem.

\begin{assumption}[Ordinary dielectric class]\label{pr:class}
Fix $\omega>0$, $k_0=\omega\sqrt{\eps_0\mu_0}$, and $1<\eps_+<\infty$. Use $e^{-i\omega t}$ and
\begin{equation}\label{pr:maxwell}
 \nabla\times E=i\omega\mu_0H,\qquad
 \nabla\times H=-i\omega\eps_0\eps_r E,
 \qquad 1\leq\eps_r(x)\leq\eps_+.
\end{equation}
The medium is real, scalar, reciprocal and nonmagnetic. There are no impressed interior sources, conductors, gain, or negative permittivities. Each geometry is piecewise smooth with finitely many pieces, but that number and all feature scales can vary. The exterior background is vacuum. Compactly supported three-dimensional dielectric perturbations form the default radiation class. Ideal leads and translationally invariant systems must instead be expressly labeled and supplied with their own exterior/port conventions.
\end{assumption}
The case $\eps_+=1$ is homogeneous and is excluded only from hotspot constructions that need contrast. A broader positive class fixes the \emph{same} $0<\eps_-\leq\eps_r\leq\eps_+$ for every member. Allowing $\eps_-\to0$ is a different target. Tensor permittivity requires separate statements. Losslessness here concerns one frequency; a nondispersive energy interpretation is an additional assumption.

\begin{assumption}[Scattering and power]\label{pr:scattering}
For each device specify a well-posed outgoing Maxwell scattering problem
with incoming and selected outgoing Hilbert spaces
$\mathcal H_{\rm in},\mathcal H_{\rm out}$. Let
$T:\mathcal H_{\rm in}\to\mathcal H_{\rm out}$ be the selected scattering
map. The physical incoming power is $\|a\|^2$ and the selected output
power is $\|Ta\|^2$. The total field depends linearly on $a$, belongs to
$H_{\rm loc}(\operatorname{curl})$, and includes all feedback. The full
scattering map includes reflected and unselected radiation. Thus
$T^*T\leq I$. At any exceptional frequency a finite, unique
limiting-absorption field must be justified separately; adding a freely
chosen trapped mode is forbidden.
\end{assumption}
This assumption is an admissibility condition, not a theorem that arbitrary rough coefficients yield a scattering theory. Radiation channels can be normalized by incoming spherical-wave flux. An ideal waveguide uses its complete propagating-mode power form. Neither convention counts an infinite-power plane wave as a unit-power input.

\begin{definition}[Counted subspaces through one physical section]\label{pr:count}
Fix the square $D=(-L/2,L/2)^2\times\{0\}$, its orientation $+\hat z$, and
$A=L^2$, together with a source-free neighborhood of $\overline D$.
Fix $0<\tau<1$, $\eta>0$ and $\delta\geq0$. For the total fields
$(E_a,H_a)$ corresponding to $a\in\mathcal H_{\rm in}$, write
$S_z(a)=\tfrac12\operatorname{Re}(E_a\times\overline{H_a})\cdot\hat z$
for the normal component of the time-averaged Poynting vector. A linear
subspace $W\subset\mathcal H_{\rm in}$ is counted if \emph{every}
$a\in W$ satisfies
\begin{align}
 \|Ta\|^2&\geq\tau\|a\|^2,\label{pr:transmission}\\
 P_D(a):=\tfrac12\operatorname{Re}\int_D(E_a\times\overline{H_a})\cdot\hat z\,\dd A
 &\geq\eta\|Ta\|^2,\label{pr:passage}\\
 \int_{\{z=0\}\setminus D}|S_z(a)|\,\dd A&\leq\delta\|a\|^2.\label{pr:bypass}
\end{align}
These integrals require legitimate physical traces with integrable flux. Otherwise a specified regularized measurement must replace them, with its own target. No product of arbitrary distributional tangential traces is assumed defined. The passage interpretation is operational; $D$ need not be a sole topological connection.
\end{definition}
Using a square freezes shape. An extension to other fixed Lipschitz sections must list shape constants. In particular one cannot replace $D$ with increasingly long thin sections of equal area while claiming a fixed-section counterexample.

\begin{conjecture}[Primary channel target: DISPROVED as quantified]\label{pr:channel-target}
For fixed parameters in Definition~\ref{pr:count} and Assumptions~\ref{pr:class}--\ref{pr:scattering}, there is a finite bound on $\dim W$ uniform over admissible devices. A stronger target for scaled copies of the same shape is
\begin{equation}\label{pr:area-bound}
 \dim W\leq C(\tau,\eta,\delta,\eps_-/\eps_+,\text{declared shape data})
                 (1+k_0^2\eps_+A).
\end{equation}
Surrounding support, resonator quality factor, matching structures, port number, interface count and minimum feature size are not bounded resources of the area-only conjecture. A theorem depending on any of them must say so.
\end{conjecture}
Finiteness precedes the scaling question; no optimal coefficient is requested. Theorem~\ref{uc:counterexample} disproves this assertion for every fixed $\delta>\eta>0$ when all outgoing radiation is selected. It does not settle a modified assertion with $\delta\leq\eta$ or specified receiving outputs. Fixing total support is also a different problem. Our flux-resource result additionally needs $\delta<\eta\tau$; other parameter values are not silently discarded from the target.

If $W_\tau=\Ran\mathbf1_{[\tau,1]}(T^*T)$ satisfies both passage conditions, a bound on counted subspaces bounds $\dim W_\tau$. Without that hypothesis there may be unrelated efficient channels bypassing $D$. Without compactness, the spectral projection is the definition; a discrete singular-value decomposition is not assumed. Nonzero coupling alone uses no positive threshold and answers a different question.

\begin{definition}[Separate hotspot targets]\label{pr:hotspot-target}
Fix contrast, frequency, a nonempty open patch $\mathcal P$, $N\in\N$, a localization length $\ell>0$ (as small relative to $\lambda_0=2\pi/k_0$ as desired), and a peak target $M>0$. The driven target seeks one admissible medium and one physical unit-power incident field with $N$ disjoint observation neighborhoods in $\mathcal P$, each containing a peak exceeding $M$ and having a controlled boundary below a fixed fraction of its peak. In a two-dimensional reduction power is per invariant length and localized regions are transverse, not finite three-dimensional points. A modal target instead requires a genuine mode and a defined modal normalization, for example the relevant electric-energy integral divided by local electric energy density. A driven monochromatic scattering field has infinite all-space energy and cannot be inserted into that formula.
\end{definition}
No $N$-independent excitation coefficient or nonvanishing fraction of total energy per hotspot is required. Such conditions would define substantially stronger problems.

\subsection{Relation to previous work}
Table~\ref{pr:prior-art} summarizes relevant results and the hypotheses
needed to compare them with the formulations above.
\begin{longtable}{@{}p{.20\textwidth}p{.38\textwidth}p{.35\textwidth}@{}}
\caption{Selected prior work and its relation to the present questions.}\label{pr:prior-art}\\
\toprule Source & Scope & Distinction from the present questions\\\midrule\endfirsthead
\toprule Source & Scope & Distinction from the present questions\\\midrule\endhead
Miller (2023), S2/S12~\cite{miller2023} & Uniform-medium channel counting motivates thickness estimates. S2 labels the nonuniform-index replacement conjectural; S12 distinguishes approximate aperture counting and design-specific Green operators. & A uniform estimate for actual structured interiors, including their total resonant response.\\[4pt]
Kuang et al. (2025), II/E~\cite{kuang2025} & Background Green operators between domains separable by a spherical surface; sizes and separations enter. Power normalization adds source/receiver response factors. & Their Eq.~(1) retains the norms of those factors. Appendix E rejects general monotonicity of singular-value ratios and a step-shaped spectrum assumption.\\[4pt]
Costabel et al. (1999)~\cite{costabel1999} & Piecewise constant positive coefficients on a bounded polyhedral partition; singular parts are linked to scalar transmission problems with stated boundary conditions. & Local singular structure alone does not prove global excitation, outgoing scattering, or simultaneous peaks.\\[4pt]
Zhao et al. (2020)~\cite{zhao2020} & Mode-volume optimization with Maxwell constraints, refractive-index limits, minimum feature constraints and numerical dual bounds; polarization matters. & Numerical feature-size scaling is neither a packing proof nor a channel theorem. Driven fields need a different normalization.\\
\bottomrule
\end{longtable}

Harrach, Lin, and Liu~\cite{harrach2018} establish field localization and
a partial-boundary-data Runge approximation result for Maxwell equations
in bounded domains. Their boundary-input setting differs from the present
open-scattering constructions and their power normalization. Bille-Lauridsen
and M\o rk~\cite{billeLauridsen2026} analyze singularity regularization in
dielectric bowtie geometries using electrostatic and quasinormal-mode
calculations; their local scaling results do not imply simultaneous global
hotspot packing.

The supplied 33-page report~\cite{priorreport} proves a full-vector count with an explicit local field-norm resource, and identifies Bragg amplification, evanescent local flux, and resonant-feedback obstructions. We audit and retain these with their restrictions. Its ENZ slot construction varies both the lower permittivity bound and the section shape; that observation alone excludes it as a counterexample to Conjecture~\ref{pr:channel-target}. We do not claim an independent recertification of every ENZ numerical result.

The supplied HTML notes~\cite{priorhtml} concern invariant slabs and fibers, scalar compact-support scattering, and finite-dimensional formal checks. Their fixed-propagation-constant magnetic variational form is monotone in $\eps^{-1}$ on a common divergence-free domain. Converting that count to a count of all propagation constants at fixed frequency needs additional spectral arguments; the notes themselves flag that step as numerically checked only. Neither a slab count nor such a fiber comparison is a full-device section theorem. The companion file named \texttt{optical\_channel\_bounds.pdf} was absent from the supplied Downloads location. No unprovided Lean or simulation archive was reconstructed or represented as verified.

A targeted search through 17 October 2026 also inspected primary recent work on metawaveguide microrings~\cite{gafsi2025}, singular dielectric cavities~\cite{mao2025}, and multi-plane communication-mode synthesis~\cite{deAngelis2026}. These concern specific confinement or wavefront-synthesis systems. None of the inspected material supplies the all-device fixed-section theorem with the quantifiers above. This is a limited search outcome, not a claim of exhaustive novelty or absence of a solution elsewhere.
